\begin{abstract}
Retrieval-augmented language models receive evidence piece by piece, and at some point that evidence becomes enough to answer the question. We ask whether the model's internal state registers this moment, and whether it also registers whether the model will actually use the evidence. We build evidence trajectories for multi-hop questions: each gold paragraph is added in turn, and every decisive addition is paired with matched additions of irrelevant, duplicate or falsified paragraphs to the same starting context. Linear probes on the residual stream at the answer position, read before any token is generated, recognise the decisive addition at $0.95$ AUROC in three model families and predict whether the model will take up the decisive evidence rather than ignore it at $0.78$. We compare these readouts with the baselines a practitioner would otherwise use. For sufficiency, the internal state is far more reliable than the model's own verbalised judgement ($0.98$ against $0.71$--$0.77$ AUROC with paragraph count matched): the models declare complete evidence insufficient a third of the time, and the probe recovers most of those cases. For uptake, text classifiers stay near chance, whether the model answers at all already predicts much of it, and the internal state adds to that signal and is the best single predictor of correctness among the answers the model gives. The probe directions turn out to be readouts rather than levers: steering along them changes nothing, whereas the difference between the mean sufficient and insufficient states acts as a gate on abstention. The readouts replicate on HotpotQA, on RAGBench's naturally retrieved documents across five domains, on a synthetic benchmark and on BM25- and E5-retrieved Wikipedia evidence, with directions learned on Wikipedia transferring only partially to other domains.
\end{abstract}
